% Suggested replacement or addition for the beginning of Section 9.
% This version avoids BibTeX dependencies; citations can be added later.

\subsection{Recommended numerical directions}

There are many possible numerical directions, but I think the two most natural
ones are:
\[
    \text{(i) censored-observation data assimilation / inverse modeling,}
    \qquad
    \text{(ii) direct surrogate learning with a censored loss.}
\]
These directions keep the project close to uncertainty quantification,
inverse problems, data assimilation, and physics-informed surrogate modeling.
They also connect naturally to the current toy experiments: the same distinction
between exact, discard, hinge, and censored treatments can be tested first in a
controlled setting and then in a more physical model.

\subsubsection{Censored-observation data assimilation / inverse modeling}

The cleanest next step is to view saturated thermal images as a data assimilation
or inverse problem with a nonstandard observation operator.  Let $T_\theta$
denote a temperature field determined by unknown parameters $\theta$, such as
heat-source parameters, material parameters, boundary flux parameters, or a
low-dimensional representation of the hidden thermal state.  In a physics-based
version, $T_\theta$ would be constrained by a heat model such as
\[
    \rho C_p \frac{\partial T_\theta}{\partial t}
    - \nabla\cdot(\kappa\nabla T_\theta)=0,
    \qquad
    -\kappa\nabla T_\theta\cdot n = h_\theta .
\]
If $\mathcal H_i(T_\theta)$ is the model-predicted temperature at pixel $i$,
the censored observation model is
\[
    Y_i = \min\{\mathcal H_i(T_\theta)+\varepsilon_i,c\},
    \qquad
    \varepsilon_i\sim\mathcal N(0,\sigma^2),
\]
where $c$ is the camera saturation threshold.  Thus the unsaturated and saturated
sets are
\[
    \mathcal U=\{i:Y_i<c\},\qquad
    \mathcal S=\{i:Y_i=c\}.
\]
A MAP-type objective for estimating $\theta$ is then
\[
    J_{\mathrm{cens}}(\theta)
    =
    \frac{1}{2\sigma^2}\sum_{i\in\mathcal U}
        \bigl(Y_i-\mathcal H_i(T_\theta)\bigr)^2
    -
    \sum_{i\in\mathcal S}
        \log \Phi\!\left(
        \frac{\mathcal H_i(T_\theta)-c}{\sigma}
        \right)
    + R(\theta),
\]
where $R(\theta)$ is a prior or regularization term.  This is a direct
physics-based analogue of the censored loss already tested in the toy Gaussian
experiment.

The main research question would be: how much do saturated pixels reduce
uncertainty in physically meaningful quantities of interest?
Examples of such quantities include peak temperature, total heat, melt-pool
width, source location, laser power, or absorptivity.  This framing is attractive
because the saturation mask is not merely missing data; it is an event
observation of the form
\[
    \mathcal H_i(T_\theta)+\varepsilon_i \geq c.
\]
That makes it closely related to feature-informed data assimilation: the
observation is not just a point value, but information that a threshold-crossing
event occurred.

A feasible study could proceed in stages:
\begin{enumerate}
    \item Start with the current Gaussian field and estimate a small number of
    parameters $\theta$.
    \item Move to synthetic fields generated by a heat-equation or thermal
    simulation model, where the true field and true parameters are known.
    \item Compare exact, discard, hinge, and censored treatments using both
    reconstruction error and uncertainty in quantities of interest.
\end{enumerate}
This direction is probably the best conceptual fit for the project because it
connects the camera-saturation problem to inverse modeling, uncertainty
quantification, and data assimilation.

\subsubsection{Direct surrogate learning with censored thermal observations}

A second promising direction is to ask whether a surrogate model can be trained
directly from censored thermal images, without first reconstructing the hidden
temperature field.  Let $a_n$ denote the input conditions for training example
$n$, such as process parameters, source parameters, or previous state
information, and let
\[
    \widehat T_\eta(a_n;x_i)
\]
be the surrogate-predicted temperature at pixel $i$.  Instead of training on the
clipped image as if it were exact, the surrogate can be trained with the censored
loss
\[
    \mathcal L_{\mathrm{cens}}(\eta)
    =
    \sum_n
    \left[
    \frac{1}{2\sigma^2}\sum_{i\in\mathcal U_n}
        \bigl(Y_{n,i}-\widehat T_\eta(a_n;x_i)\bigr)^2
    -
    \sum_{i\in\mathcal S_n}
        \log \Phi\!\left(
        \frac{\widehat T_\eta(a_n;x_i)-c_n}{\sigma}
        \right)
    \right].
\]
If desired, this can be combined with a physics-informed penalty,
\[
    \mathcal L(\eta)
    =
    \mathcal L_{\mathrm{cens}}(\eta)
    +
    \gamma\sum_n
    \left\|\mathcal F\bigl(\widehat T_\eta(a_n)\bigr)\right\|^2,
\]
where $\mathcal F$ is a heat-equation residual or another physics-based
constraint.

The empirical question would be: is it better to train a surrogate directly from
censored data, or to reconstruct the full field first and then train on the
reconstructions?
This can be tested cleanly with synthetic data, where the true unclipped field is
available.  A useful comparison would include:
\[
\begin{array}{ll}
    \text{clipped training} & \text{treat } Y_i \text{ as exact everywhere},\\
    \text{discard training} & \text{train only on } i\in\mathcal U,\\
    \text{censored training} & \text{use the likelihood above},\\
    \text{oracle training} & \text{train on the true unclipped field.}
\end{array}
\]
This direction is more machine-learning oriented than the first one, but it is a
natural extension if the long-term goal is to build useful surrogate models from
thermal-camera data.  I would treat fully generative reconstruction methods as a
later extension, after the censored observation model and its value for
uncertainty reduction are clear.
